\chapter{\eh{truck} Capa 4 --- Transporte (Transport Layer)}

\section{\eh{thinking-face} ¿Qué hace la capa de transporte?}

\begin{concepto}
La capa de transporte entrega datos de \textbf{proceso a proceso}
(aplicación a aplicación) usando \textbf{números de puerto}.
Vive en los hosts extremos --- los routers intermedios no la procesan.

\emoji{light-bulb} \textbf{Analogía:} IP sabe llegar a la \textit{casa} correcta (host).
El puerto es el \textit{número de departamento} en esa casa. El cartero (IP)
deja el paquete en la puerta del edificio; el portero (transporte) lo sube
al departamento correcto (la aplicación).
\end{concepto}

\subsection{Números de puerto}

\begin{table}[h!]
\centering
\renewcommand{\arraystretch}{1.3}
\begin{tabular}{llll}
\toprule
\textbf{Rango} & \textbf{Tipo} & \textbf{Ejemplos} & \textbf{Quién los asigna} \\
\midrule
0 -- 1023    & Well-known    & 80=HTTP, 443=HTTPS, 22=SSH, 53=DNS & IANA \\
1024 -- 49151 & Registered   & 3306=MySQL, 5432=PostgreSQL, 6379=Redis & IANA \\
49152 -- 65535 & Ephemeral   & Origen temporal del cliente & OS \\
\bottomrule
\end{tabular}
\caption{Rangos de números de puerto}
\end{table}

Una \textbf{socket} = IP + puerto. Identifica únicamente un proceso en la red:
\texttt{192.168.1.10:54321} ↔ \texttt{142.250.80.46:443}

\section{\eh{zap} UDP --- User Datagram Protocol (RFC 768)}

\begin{concepto}
\textbf{UDP} es \textbf{connectionless} (sin conexión previa): envía datagramas
sin establecer conexión, sin confirmar que llegaron, sin controlar el orden.

\emoji{light-bulb} \textbf{Analogía:} UDP es como mandar postales.
Las mandas y asumes que llegan. Si se pierde una... ni modo.
Pero es \textit{mucho más rápido} que mandar carta certificada.
\end{concepto}

Header UDP --- solo \textbf{8 bytes}:

\begin{lstlisting}
  0              15 16             31
  ┌───────────────┬─────────────────┐
  │  Source Port  │ Destination Port│
  │   (16 bits)   │   (16 bits)     │
  ├───────────────┼─────────────────┤
  │    Length     │    Checksum     │
  │   (16 bits)   │   (16 bits)     │
  └───────────────┴─────────────────┘
\end{lstlisting}

Casos de uso de UDP: DNS, DHCP, TFTP, streaming de video (YouTube usa QUIC/UDP),
VoIP, gaming online, NTP. \textbf{Velocidad $>$ confiabilidad.}

\section{\eh{handshake} TCP --- Transmission Control Protocol (RFC 793)}

\begin{concepto}
\textbf{TCP} es \textbf{connection-oriented} (orientado a conexión):
establece una conexión antes de enviar datos, garantiza \textbf{entrega en orden},
retransmite paquetes perdidos y controla el flujo para no saturar al receptor.

\emoji{light-bulb} \textbf{Analogía:} TCP es como una llamada telefónica.
Primero marcas (SYN), esperas que respondan (SYN-ACK), confirmas
``¡Hola, te escucho!'' (ACK), y \textit{después} empiezas a hablar.
Además, si no escuchaste una parte, dices ``¿puedes repetir?'' y la otra
persona lo repite (retransmisión).
\end{concepto}

\subsection{Header TCP --- 20 bytes mínimo}

\begin{lstlisting}
  0         1         2         3
  ┌──────────────────┬──────────────────┐
  │   Source Port    │  Dest. Port      │  (32 bits)
  ├──────────────────┴──────────────────┤
  │           Sequence Number           │  (32 bits)
  ├─────────────────────────────────────┤
  │         Acknowledgment Number       │  (32 bits)
  ├────┬────┬─────────────┬─────────────┤
  │Off │Rsv │   Flags     │   Window    │  (32 bits)
  │4b  │4b  │U A P R S F  │  (16 bits)  │
  ├────┴────┴─────────────┴─────────────┤
  │    Checksum         │  Urgent Ptr   │  (32 bits)
  └─────────────────────────────────────┘
\end{lstlisting}

\subsection{Flags TCP}

\begin{table}[h!]
\centering
\renewcommand{\arraystretch}{1.3}
\begin{tabular}{clll}
\toprule
\textbf{Bit} & \textbf{Flag} & \textbf{Nombre} & \textbf{Uso} \\
\midrule
1 & URG & Urgent    & Datos urgentes, procesarlos antes \\
2 & ACK & Acknowledge & Confirma recepción (Ack Number es válido) \\
3 & PSH & Push      & Enviar al app inmediatamente, sin buffering \\
4 & RST & Reset     & Cierre abrupto de la conexión / error \\
5 & SYN & Synchronize & Iniciar conexión (intercambiar Seq Numbers) \\
6 & FIN & Finish    & Cierre normal de la conexión \\
\bottomrule
\end{tabular}
\caption{Flags TCP}
\end{table}

\subsection{TCP 3-Way Handshake}

\begin{center}
\begin{tikzpicture}[
    msg/.style={-{Stealth}, thick},
    nota/.style={font=\small\ttfamily, text=gray!70!white}
]
\node[font=\bfseries, text=textlight] (client) at (0,  0) {Cliente};
\node[font=\bfseries, text=textlight] (server) at (8,  0) {Servidor};
\draw[thick, draw=gray!50!white] (0,-0.3) -- (0,-5.5);
\draw[thick, draw=gray!50!white] (8,-0.3) -- (8,-5.5);

% SYN
\draw[msg, cyan!70!white] (0,-1.0) -- (8,-1.8)
    node[midway, above, sloped, font=\small, text=cyan!80!white] {\textbf{SYN} seq=1000};
\node[nota, left] at (0,-1.0) {CLOSED};
\node[nota, right] at (8,-1.8) {SYN\_RCVD};

% SYN-ACK
\draw[msg, orange!80!white] (8,-1.8) -- (0,-2.8)
    node[midway, above, sloped, font=\small, text=orange!80!white] {\textbf{SYN+ACK} seq=5000, ack=1001};
\node[nota, left] at (0,-2.8) {SYN\_SENT};

% ACK
\draw[msg, green!60!white] (0,-2.8) -- (8,-3.6)
    node[midway, above, sloped, font=\small] {\textbf{ACK} ack=5001};
\node[nota, right] at (8,-3.6) {ESTABLISHED};
\node[nota, left] at (0,-3.6) {ESTABLISHED};

\node[font=\small\itshape, text=green!60!white] at (4,-4.4)
    {Conexión establecida \emoji{check-mark-button} --- datos pueden fluir};
\end{tikzpicture}
\end{center}

\begin{cuidado}
\textbf{SYN Flood Attack:} un atacante envía miles de SYNs sin completar
el handshake (nunca manda el ACK final). El servidor guarda estado en
memoria por cada SYN → agota recursos (DoS).

Defensa: \textbf{SYN cookies} (el servidor no guarda estado hasta recibir
el ACK), límite de tasa de SYNs, firewalls.
\end{cuidado}

\subsection{TCP 4-Way Termination}

\begin{lstlisting}
  Cliente              Servidor
     │─── FIN ──────────►│   "Terminé de enviar datos"
     │◄── ACK ───────────│   "Recibido, espera"
     │                   │   (servidor puede seguir enviando)
     │◄── FIN ───────────│   "Yo también terminé"
     │─── ACK ──────────►│   "Ok, cerramos"
     │                   │
    [TIME_WAIT: 2×MSL ≈ 4 min]
\end{lstlisting}

\textbf{TIME\_WAIT:} el cliente espera para asegurarse de que el último ACK
llegó. Si el servidor no lo recibió, retransmitirá el FIN y el cliente podrá
responder. Esto previene que segmentos ``fantasmas'' de la conexión anterior
contaminen una nueva.

\subsection{Control de flujo y congestión}

\begin{concepto}
\textbf{Control de flujo:} el receptor anuncia su \textbf{receive window (rwnd)}
en cada ACK: ``puedo recibir hasta X bytes más sin esperar''. El emisor
no envía más de lo que el receptor puede procesar.

\textbf{Control de congestión:} TCP ajusta la velocidad según el estado
de la red, controlado por la \textbf{congestion window (cwnd)}.
\end{concepto}

Fases del control de congestión TCP (algoritmo CUBIC, default en Linux):

\begin{enumerate}
    \item \textbf{Slow Start:} cwnd comienza en 1 MSS y se \textit{dobla}
          con cada RTT. Crecimiento exponencial hasta llegar a \textit{ssthresh}.
    \item \textbf{Congestion Avoidance:} al superar ssthresh, crece en
          1 MSS por RTT (lineal). Más conservador.
    \item \textbf{Fast Retransmit:} si llegan \textbf{3 ACKs duplicados},
          el emisor retransmite inmediatamente sin esperar timeout.
    \item \textbf{Fast Recovery:} tras la pérdida detectada por 3-dup-ACK,
          ssthresh = cwnd/2 y cwnd = ssthresh (no regresa a 1). Rápida recuperación.
    \item \textbf{Timeout:} si no llega ACK en el timeout, cwnd regresa a 1.
          Mucho más severo que fast recovery.
\end{enumerate}

\begin{cuidado}
\textbf{¿Por qué TCP tiene problemas con HTTP?}

Cada nueva conexión TCP debe hacer handshake (1 RTT) + slow start desde cero.
HTTP/1.1 abría múltiples conexiones TCP en paralelo (hasta 6 por dominio)
para simular paralelismo, pero esto saturaba la red.

HTTP/2 resolvió esto con \textit{multiplexing} sobre una sola conexión TCP.
Pero TCP en sí tiene un problema: si un paquete se pierde, todos los streams
esperan (Head-of-Line Blocking a nivel TCP). QUIC (ver capítulo 8)
resuelve esto usando UDP con streams independientes.
\end{cuidado}
